\section{Projective surfaces and integral blow-ups}
\label{sec:geometry}

We need two projective surfaces whose integral homology distinguishes
the primes $2$ and $3$. Their products with projective spaces will serve
as blow-up centers. The integral blow-up formula will then place their
torsion in a simply connected K\"ahler manifold without introducing any
extension ambiguity.

\subsection{The two surfaces}

Let
\[
 E=\C/(\Z+i\Z),\qquad
 \zeta_p=\exp(2\pi i/p),\qquad
 \Lambda_2=\Z+i\Z,\qquad
 \Lambda_3=\Z+\zeta_3\Z,
\]
and put $E_p=\C/\Lambda_p$ for $p=2,3$. The automorphism
\[
 \sigma_p:E\times E_p\longrightarrow E\times E_p,
 \qquad (z,w)\longmapsto(z+1/p,\zeta_p w)
\]
has order $p$. Every nonidentity power acts freely, because its
translation on the first factor is nontrivial. Thus the quotient
\[
 S_p=(E\times E_p)/\langle\sigma_p\rangle
\]
is a smooth compact connected complex surface.
These are bielliptic surfaces; see \cite[Section~2]{Blomme2025} for
their place in the classification. The explicit embeddings and
computations below are those of
\cite[Lemma~3.1]{OpenAI2026Stable}.

\begin{lemma}
\label{geo:surfaces}
The surfaces $S_2$ and $S_3$ admit holomorphic embeddings into
$\CP^{44}$ and $\CP^{164}$, respectively. If
$T_2=(\Z/2)^2$ and $T_3=\Z/3$, then
\begin{equation}
\label{geo:surface-homology}
 H_i(S_p;\Z)\cong
 \begin{cases}
  \Z,&i=0,4,\\
  \Z^2\oplus T_p,&i=1,2,\\
  \Z^2,&i=3,\\
  0,&\text{otherwise}.
 \end{cases}
\end{equation}
In particular, $b(S_p)=8$.
\end{lemma}

\begin{proof}
We first give an embedding with an explicit ambient dimension. A
degree-three line bundle on an elliptic curve has a three-dimensional
space of sections. Riemann--Roch gives dimension two or one after
subtracting an effective divisor of degree one or two. Its complete
linear system consequently has no base points and separates points
and tangent vectors. It embeds
the curve into $\CP^2$. Applying this to both factors and taking the
Segre embedding yields
\[
 \iota_p:E\times E_p\hookrightarrow\CP^8.
\]

Regard a nonzero vector in $\C^9$ as a linear polynomial on its dual.
Multiplying representatives of the $p$ points
\[
 \iota_p(x),\ \iota_p(\sigma_p x),\ \ldots,
 \iota_p(\sigma_p^{p-1}x)
\]
defines a projective degree-$p$ polynomial. This construction is
holomorphic and independent of the representatives. It is unchanged
by $\sigma_p$, so it descends to a holomorphic map
\[
 F_p:S_p\longrightarrow\mathbb P(\operatorname{Sym}^p\C^9).
\]
Unique factorization recovers the unordered orbit from its product of
linear factors; hence $F_p$ is injective.

To check immersion, choose local nonzero representatives $L_j$ of the
orbit points. If a tangent vector has zero image under the projective
differential of $F_p$, its derivatives satisfy
\[
 \sum_{j=0}^{p-1}\dot L_j\prod_{k\ne j}L_k
   =a\prod_{k=0}^{p-1}L_k
\]
for a scalar $a$. Freeness makes the $L_j$ pairwise nonproportional.
Restricting this identity to $\{L_j=0\}$ gives $\dot L_j=0$ there:
the other factors have nonzero product in the polynomial ring of that
hyperplane, which is an integral domain. Thus each $\dot L_j$ is
proportional to $L_j$. Since $\iota_p$ is immersive and the quotient
map is a local biholomorphism, the original tangent vector vanishes.
Compactness now makes $F_p$ a holomorphic embedding. The dimensions
\[
 \dim\operatorname{Sym}^2\C^9=45,
 \qquad \dim\operatorname{Sym}^3\C^9=165
\]
give the stated projective spaces.

For the integral homology, use the universal covering $\C^2\to S_p$.
Its deck group is generated by
\[
 \begin{aligned}
 a(z,w)&=(z+1/p,\zeta_p w),&d(z,w)&=(z+i,w),\\
 t_\lambda(z,w)&=(z,w+\lambda),&&\lambda\in\Lambda_p.
 \end{aligned}
\]
Indeed, $a^p$ is translation by $1$ in the first coordinate, and these
generators therefore include the full lattice of $E\times E_p$ and
a lift of $\sigma_p$. The commuting transformations $a,d$ generate
$\Z^2$, act on $\Lambda_p$ by $\zeta_p$ and the identity, and have
trivial intersection with its translation subgroup. Consequently
\[
 \begin{aligned}
 \pi_1(S_p)&\cong\Lambda_p\rtimes\Z^2,\\
 H_1(S_p;\Z)&\cong
 \Z^2\oplus\Lambda_p/(\zeta_p-1)\Lambda_p.
 \end{aligned}
\]
The lattice endomorphism $\zeta_p-1$ is $-2I$ for $p=2$. For $p=3$,
its matrix in the basis $(1,\zeta_3)$ is
\[
 \begin{pmatrix}-1&-1\\1&-2\end{pmatrix},
\]
whose Smith normal form is $\operatorname{diag}(1,3)$. This proves the
asserted description of $H_1$.

The complex orientation gives integral Poincar\'e duality. In
particular, $H_3(S_p;\Z)\cong H^1(S_p;\Z)\cong\Z^2$, while the
universal coefficient theorem gives
\[
 \Tor H_2(S_p;\Z)\cong\Tor H^2(S_p;\Z)
 \cong\Tor H_1(S_p;\Z)\cong T_p.
\]
Finally, the degree-$p$ covering by $E\times E_p$ implies
$p\chi(S_p)=\chi(E\times E_p)=0$. The ranks already computed force
$\rank H_2(S_p;\Z)=2$, completing
\eqref{geo:surface-homology}.
\end{proof}

\subsection{The integral blow-up formula}

The following standard formula is integral, not merely rational; see
\cite[Theorem~4.1]{DuanLi2009}. We include the splitting argument to
make explicit why the torsion in Lemma~\ref{geo:surfaces} survives.

\begin{proposition}
\label{geo:blowup}
Let $A$ be a closed connected K\"ahler manifold and let $Z\subset A$
be a closed connected complex submanifold of complex codimension
$c\geq2$. The complex blow-up $\pi:\widehat A\to A$ is a closed
connected K\"ahler manifold, and it is simply connected if $A$ is.
For every integer $i$,
\begin{equation}
\label{geo:blowup-homology}
 H_i(\widehat A;\Z)\cong
 H_i(A;\Z)\oplus\bigoplus_{j=1}^{c-1}H_{i-2j}(Z;\Z).
\end{equation}
\end{proposition}

\begin{proof}
Write $\nu$ for the complex normal bundle of $Z$ in $A$. The local
blow-up model replaces $v\in\C^c$ by pairs $(\ell,v)$ with
$v\in\ell$. It gives a smooth complex manifold, a proper map $\pi$
that is an isomorphism off $Z$, and an exceptional divisor
\[
 E=\mathbb P(\nu),\qquad \rho:E\to Z.
\]
The normal line bundle of $E$ is $\mathcal O_E(-1)$. Properness and
the local model give compactness and connectedness.

Choose a K\"ahler form $\omega_A$ and a Hermitian metric on $\nu$.
The induced metric on $\mathcal O_E(1)$ has curvature positive along
the projective fibers. Since
$\mathcal O(-E)|_E\cong\mathcal O_E(1)$, extend this metric to
$\mathcal O(-E)$ and let $\eta$ be its real Chern curvature form.
Such an extension is obtained by extending the logarithm of the ratio
with any ambient metric. Curvature commutes with restriction to the
complex submanifold $E$, so $\eta$ retains fiberwise positivity.
The form $\pi^*\omega_A$ is semipositive, with kernel along $E$
equal to the tangent space of the fiber of $\rho$, and is positive
off $E$. On these kernel directions $\eta$ is positive. Compactness
of the unit tangent bundle therefore implies that
\[
 \pi^*\omega_A+\varepsilon\eta
\]
is positive for all sufficiently small $\varepsilon>0$: near the
kernel directions $\eta$ is positive, and on their complement
$\pi^*\omega_A$ has a positive lower bound whereas $\eta$ is bounded.
This closed real $(1,1)$-form is K\"ahler.

If $A$ is simply connected, then so is $A\setminus Z$. Indeed,
general position allows a contracting disk for any loop in that
complement to avoid $Z$, whose real codimension is at least four.
A loop in $\widehat A$ can similarly be perturbed off the real
codimension-two divisor $E$. It then contracts in
$\widehat A\setminus E\cong A\setminus Z$.

For the remaining assertion, all cohomology groups have coefficients
in $\Z$. Let $\iota:E\hookrightarrow\widehat A$ and
$h=c_1(\mathcal O_E(1))$. The integral Leray--Hirsch theorem gives
\begin{equation}
\label{geo:projective-bundle}
 H^k(E)=\bigoplus_{j=0}^{c-1}h^j\rho^*H^{k-2j}(Z).
\end{equation}
This uses the integral fiber basis $1,h,\ldots,h^{c-1}$ and does not
require torsion-free cohomology of $Z$. The integral Gysin push-forward
$\iota_*:H^k(E)\to H^{k+2}(\widehat A)$ satisfies
$\iota^*\iota_*(v)=-hv$, because the normal bundle of $E$ is
$\mathcal O_E(-1)$. Thus the classes
\[
 s_j(u)=-\iota_*(h^{j-1}\rho^*u),
 \qquad u\in H^{k-2j}(Z),\quad 1\leq j\leq c-1,
\]
restrict to $h^j\rho^*u$ on $E$.

We first show that a class $v\in H^k(\widehat A)$ whose restriction
to $E$ is $\rho^*z$ lies in $\pi^*H^k(A)$. Collapsing the exceptional
divisor and the center gives a homeomorphism
$\widehat A/E\cong A/Z$, so $\pi$ induces an isomorphism
$H^*(A,Z)\cong H^*(\widehat A,E)$. In the long exact sequences of
these pairs, the connecting image of $\rho^*z=\iota^*v$ is zero.
Naturality and this relative isomorphism show that the connecting
image of $z$ is zero as well. Hence $z$ is the restriction of some
$u\in H^k(A)$. The class $v-\pi^*u$ vanishes on $E$, so it comes
from relative cohomology; the same relative isomorphism expresses
it as the pullback of a class on $A$. This proves the claim.

Given an arbitrary $v$, subtracting suitable $s_j(u_j)$ removes the
positive powers of $h$ in \eqref{geo:projective-bundle}. The claim
then proves surjectivity of
\begin{equation}
\label{geo:blowup-cohomology-map}
 \pi^*+\sum_{j=1}^{c-1}s_j:
 H^k(A)\oplus\bigoplus_{j=1}^{c-1}H^{k-2j}(Z)
 \longrightarrow H^k(\widehat A).
\end{equation}
For injectivity, restriction to $E$ first forces every positive-power
coefficient $u_j$ to vanish by \eqref{geo:projective-bundle}.
Moreover, $\pi^*$ is injective over $\Z$. The map $\pi$ has degree
one, so cap-product naturality gives, for $\alpha\in H^k(A;\Z)$,
\[
 \pi_*\bigl(\pi^*\alpha\frown[\widehat A]\bigr)
 =\alpha\frown\pi_*[\widehat A]=\alpha\frown[A].
\]
Integral Poincar\'e duality is an isomorphism on all classes, including
torsion, so $\pi^*\alpha=0$ implies $\alpha=0$. Thus
\eqref{geo:blowup-cohomology-map} is an isomorphism.

If $\dim_{\C}A=m$, apply integral Poincar\'e duality in degree
$2m-i$. The center summands become $H_{i-2c+2j}(Z;\Z)$; replacing
$j$ by $c-j$ gives \eqref{geo:blowup-homology}.
\end{proof}
